\chapter{Literature Review}

Real-time financial information systems have evolved from periodic updates to high-frequency streaming platforms where milliseconds determine market advantage. Early systems relied on polling-based HTTP mechanisms, which often resulted in "stale" data that did not reflect true market conditions.

The transition to the WebSocket protocol has been foundational for the fintech industry, enabling persistent, full-duplex communication over single TCP connections \cite{websocket_rfc}. Unlike traditional RESTful polling, WebSockets allow financial servers to push price ticks and news alerts directly to clients, significantly reducing latency and bandwidth consumption in high-volatility scenarios.

Modern financial data aggregation also relies on robust Application Programming Interfaces (APIs). While simulated data provides a baseline, the integration of real-world feeds like the Yahoo Finance API allows for higher fidelity in testing and performance evaluation of low-latency systems. Fielding’s work on RESTful principles remains the standard for fetching structured historical data and news articles efficiently \cite{fielding}. However, the hybrid approach—using REST for management and WebSockets for live streams—is widely recognized as the most effective architecture for comprehensive financial dashboards.

Security and reliability are paramount in high-stakes financial systems. This project evaluates the synergy between application-layer protocols (HTTP/WS), transport reliability (TCP), and security (TLS) within a real-time context, providing a scalable model for live market data processing utilizing professional-grade Solarized Light visualization for human operators.